Los módulos presentados en esta propuesta explotan al máximo los datos disponibles. Sin embargo, existen preguntas relevantes para el negocio que los datos observacionales no pueden responder de forma definitiva, incluso con diseños econométricos sofisticados. Cuando la variación exógena es insuficiente, los métodos cuasi-experimentales enfrentan límites inherentes para identificar efectos causales.

Una alternativa es generar esa variación de manera deliberada mediante experimentos controlados aleatorizados (conocidos como RCT por sus siglas en inglés). En este tipo de diseños, la asignación del tratamiento se realiza de forma aleatoria, lo que permite aislar su efecto causal sin depender de supuestos fuertes sobre el proceso de selección o la dinámica de los datos.

En el contexto de YPF, esto podría traducirse, por ejemplo, en asignar aleatoriamente beneficios del programa Serviclub a distintos subgrupos de socios, implementar campañas de fidelización en estaciones seleccionadas al azar o introducir variaciones del formato Full en un conjunto de estaciones piloto. Este tipo de intervenciones permite medir con precisión el impacto de cada iniciativa sobre variables de interés como ventas, frecuencia de visita o consumo promedio.

El diseño de estos experimentos requiere planificación previa: definir con claridad las preguntas a responder, dimensionar adecuadamente los grupos de tratamiento y control, y coordinar la implementación con las áreas operativas. Por este motivo, una extensión natural de este proyecto sería dedicar una etapa específica a identificar aquellas preguntas prioritarias que no pueden responderse con los datos existentes y diseñar los experimentos necesarios para abordarlas.

El resultado de esa etapa sería un protocolo experimental listo para implementar, que permitiría transformar preguntas abiertas del negocio en evidencia causal sólida.